% ==================================================================
% id:       comb-024
% topic:    קומבינטוריקה
% subtopic: הסברים קומבינטוריים לזהויות
% solution: full
% ==================================================================
יהיו $0\le k\le r\le n$. בשאלה נוכיח את הזהות
\[
\binom nr\binom rk=\binom nk\binom{n-k}{r-k}.
\]
\begin{enumerate}
\item תנו הסבר קומבינטורי לזהות.
\hint{מתוך $n$ סטודנטים בוחרים ועדה של $r$ חברים, ומתוך הוועדה בוחרים ועדת משנה של $k$ חברים. בכמה דרכים אפשר לעשות את זה, אם מתחילים מהוועדה? ואם מתחילים מוועדת המשנה?}
\begin{solbox}
נספור בשתי דרכים את מספר האפשרויות לבחור, מתוך $n$ סטודנטים, ועדה של $r$ חברים וגם ועדת משנה של $k$ חברים מתוך הוועדה.

\textbf{קודם הוועדה.} את הוועדה בוחרים ב-$\binom nr$ דרכים, ואז את ועדת המשנה מתוך $r$ חברי הוועדה ב-$\binom rk$ דרכים. לפי עקרון הכפל: $\binom nr\binom rk$.

\textbf{קודם ועדת המשנה.} את ועדת המשנה בוחרים ישר מתוך כל הסטודנטים, ב-$\binom nk$ דרכים. אחר כך משלימים את הוועדה: צריך עוד $r-k$ חברים מתוך $n-k$ הסטודנטים שלא נבחרו, ב-$\binom{n-k}{r-k}$ דרכים. לפי עקרון הכפל: $\binom nk\binom{n-k}{r-k}$.

ספרנו את אותו הדבר בשתי דרכים, ולכן שני הביטויים שווים.
\end{solbox}
\item הוכיחו את הזהות בעזרת הנוסחה $\binom nk=\frac{n!}{k!\,(n-k)!}$.
\begin{solbox}
\[
\binom nr\binom rk=\frac{n!}{r!\,(n-r)!}\cdot\frac{r!}{k!\,(r-k)!}=\frac{n!}{(n-r)!\,k!\,(r-k)!},
\]
\[
\binom nk\binom{n-k}{r-k}=\frac{n!}{k!\,(n-k)!}\cdot\frac{(n-k)!}{(r-k)!\,(n-r)!}=\frac{n!}{k!\,(r-k)!\,(n-r)!}.
\]
(כי $(n-k)-(r-k)=n-r$.) קיבלנו את אותו ביטוי.
\end{solbox}
\end{enumerate}
